\section{Infrastructure}

\subsection{VM Deployment (Baseline)}

\begin{itemize}
  \item UTM Ubuntu 22.04 LTS virtual machine provisioned via Ansible: 4 vCPU, 8 GB RAM, 40 GB disk
  \item \dagster{} 1.12.7 deployed with \textbf{DockerRunLauncher} (each job = one Docker container on the VM host)
  \item Docker CE installed on the VM via Ansible \texttt{docker} role
  \item PostgreSQL 16 as the \dagster{} event log and run storage backend
  \item Per-container CPU/memory limits via Docker --- jobs isolated from each other but sharing the host daemon
  \item No autoscaling capability
\end{itemize}

\subsection{Kubernetes Deployment (Comparison)}

\begin{itemize}
  \item Kind (Kubernetes in Docker) single-node cluster on the same host
  \item \dagster{} 1.12.7 configured with Kubernetes Run Launcher
  \item Per-pod resource requests and limits defined per workload level
  \item Same host machine as the \vm{}, ensuring equivalent underlying hardware
  \item Fixed node capacity --- no node-level autoscaling
\end{itemize}

\subsection{Why Kind Instead of a Cloud Cluster}

Using Kind on the same host as the UTM \vm{} eliminates hardware variance as a
confound. Both environments compete for the same physical \textsc{cpu} and memory, which
means the measured differences arise from the execution model (Docker containers on a
single-host daemon vs.\ Kubernetes pods with per-pod resource limits) rather than from
infrastructure quality differences. This is a deliberate methodological choice.

\subsection{Tooling and Reproducibility}

All experimental scripts are written in Python~3.12 using \texttt{psutil} for resource
monitoring and \texttt{kubectl top} for Kubernetes metrics. \texttt{SciPy} is used for
statistical analysis. The full experimental setup is version-controlled and executable via
a single \texttt{make} target.
